\subsection{Related Works}

\subsubsection{Model-based Reinforcement Learning}

Model-based reinforcement learning approaches use a model of the environment to simulate agent trajectories, enabling better generalization, sample efficiency, and decision-making through planning. Following the success of deep neural networks for learning function approximations~\citep{lecun2015deep}, researchers proposed to learn world models using gradient backpropagation. One of the first approaches to apply model-based reinforcement learning to high-dimensional visual observations is SimPLe~\citep{kaiser2019model}. The authors proposed to learn a predictive world model for Atari games directly in pixel space using a convolutional autoencoder. SimPLe's model predicts future observations and rewards conditioned on past frames and actions, enabling the training of a Proximal Policy Optimization (PPO) agent~\citep{schulman2017proximal} using simulated rollouts.

PlaNet~\citep{hafner2019learning} subsequently introduced Recurrent State-Space Models (RSSMs), combining deterministic recurrent dynamics implemented with Gated Recurrent Units (GRUs)~\citep{cho2014learning} and stochastic latent variables to learn compact world models in latent space. The agent learns hidden representations by encoding observations using a convolutional variational autoencoder (VAE)~\citep{kingma2014auto} while predicting future latent states and rewards, enabling planning via model predictive control. Building on this framework, Dreamer~\citep{hafner2019dream} improved performance by learning a policy and value function directly within the learned latent space. DreamerV2~\citep{hafner2020mastering} extended this approach to Atari games by introducing categorical latent variables trained with straight-through gradients~\citep{bengio2013estimating}, improving representational capacity and stability in discrete domains. More recently, DreamerV3~\citep{hafner2023mastering} demonstrated that a single architecture with shared hyperparameters can scale across diverse domains, from continuous control to Atari. This was achieved through a set of architectural and optimization refinements, including symlog-scaled reward and value predictions to handle scale variation, layer normalization~\citep{ba2016layer} for improved robustness, and return normalization using an exponential moving average (EMA) of return percentiles. With these changes, DreamerV3 achieved state-of-the-art performance across a wide range of benchmarks, often outperforming both specialized model-free and other model-based methods.

In parallel to the Dreamer line of work, other approaches have explored alternative world model designs. Beyond recurrent architectures, Transformer-based world models have been proposed to improve long-term memory and representation learning~\citep{micheli2022transformers, chen2021transdreamer, robine2023transformer, zhang2023storm}. Diffusion-based world models~\citep{sohl2015deep} have also emerged as a compelling alternative to discrete latent variable sequence modeling, demonstrating strong performance across a range of settings~\citep{alonso2024diffusion, valevski2025diffusion, bar2025navigation, hafner2025training, hu2026multiplayer}. Another line of work has investigated multitask world models, training a single world model jointly across multiple tasks to improve generalization and achieve stronger performance than single-task models~\citep{hansen2024td, hansen2026learning}. Model-based reinforcement learning approaches using Monte-Carlo tree search (MCTS)~\citep{coulom2006efficient} have also been explored following the success of Alpha Go~\citep{silver2017mastering, silver2018general, schrittwieser2020mastering, ye2021mastering, wang2024efficientzero}.

\subsubsection{Memory in World Models}

Multiple approaches have focused on improving the memory capabilities of world models. \citet{pasukonis2023evaluating} proposed to augment Dreamer with Truncated Backpropagation Through Time (TBTT), conditioning training sequences on previous context representations to preserve information over time. TransDreamer~\citep{chen2021transdreamer} replaced Dreamer’s RSSM by a Transformer State-Space Model (TSSM) using masked self-attention to process past context more easily. The agent was evaluated on Hidden Order Discovery tasks requiring long-term memory and reasoning, demonstrating better performance. Recall To Imagine (R2I)~\citep{samsami2024mastering} proposed the Structured State-Space Model (S3M) to improve computational efficiency when handling long sequences. This allows R2I to process very long sequences and achieve strong performance on the Memory Maze benchmark. Concurrently,~\citet{deng2023facing} proposed S4WM, demonstrating that the S4 architecture outperforms Transformer-based
world models in terms of long-term memory. Alternatively, EDELINE~\citep{lee2025edeline} proposed to augment the DIAMOND world model~\citep{alonso2024diffusion} with a Mamba~\citep{gu2023mamba} SSM to enhance memory capabilities. A similar approach was concurrently proposed in~\citet{savov2025statespacediffuser} where the authors use Mamba to improve model predictions on CSGO~\citep{pearce2022counter}.

\subsubsection{Test-Time Training}

Test-Time Training (TTT) was originally proposed in~\citet{sun2020test} for improving the performance of predictive models when training and test data come from different distributions. This work proposed to turn a single unlabeled test sample into a self-supervised learning problem, on which to update the model parameters before making a prediction. The initial TTT paper used image rotation prediction~\citep{gidaris2018unsupervised} as a self-supervised task to update model weights and achieve higher classification accuracy on corrupted test images. The method was later improved with TTT-MAE~\citep{gandelsman2022test}, using MAE~\citep{he2022masked} instead of rotation prediction to reach higher performance. Test-time training was subsequently applied to multiple domains like reinforcement learning~\citep{hansen2021self}, language modeling~\citep{hardt2023test}, and video processing~\citep{wang2025test}. In the area of reinforcement learning, \citet{hansen2021self} proposed Policy Adaptation during Deployment (PAD)~\citep{hansen2021self}, applying TTT to model-free reinforcement learning algorithms~\citep{haarnoja2018soft, mnih2016asynchronous}. They demonstrated that PAD improves performance under distribution shifts on the DeepMind Control~\citep{tassa2018deepmind} and CRLMaze~\citep{lomonaco2020continual} benchmarks. They also showed that PAD enables more effective sim-to-real transfer for robotic manipulation tasks.